\subsection{Cancelling a quadratic pole family before integration}
\label{sec:comp-real-kernel-collection}
The separate primitive constructions above also provide a way to check
individual sectors. For the total, a shorter route combines their inner
integrals before constructing the outer primitives. The four contributions
use the same coordinate $u$, the same interval $\sqrt{s}<u<s$, and the
same relation $z=s+1-u-s/u$ to the original endpoint. They can therefore
be added at any fixed positive cutoff in $z$. This observation permits
cancellations before either a generalized polylogarithm or an endpoint
constant is evaluated.

For $A=R_0,D$, denote the rational coefficients of the corresponding
inner integral, including the Jacobian $1-s/u^2$, by $K_{A,j}$, with
$j=0,1,2,3$. Thus its outer integrand is
$K_{A,0}+\sum_{j=1}^3K_{A,j}L_j$, with the real logarithms $L_j$
already defined above. Write the energy-logarithm sector as
$\ell(K_{\log,0}+K_{\log,1}L_1+K_{\log,2}L_2)$, where
$\ell=\ln s$, and the single-$W$ sector as $A_WL_1+B_WL_1L_2$.
Here $A_W$ and $B_W$ are its two rational coefficients after the same
change of variable. The combined integrand is exactly
\begin{align}
 f_4(u)&=\overline K_0+\overline K_1L_1+
             \overline K_2L_2+\overline K_3L_3+B_WL_1L_2,\\
 \overline K_0&=K_{R_0,0}+K_{D,0}+\ell K_{\log,0},\\
 \overline K_1&=K_{R_0,1}+K_{D,1}+\ell K_{\log,1}+A_W,\\
 \overline K_2&=K_{R_0,2}+K_{D,2}+\ell K_{\log,2},\\
 \overline K_3&=K_{R_0,3}+K_{D,3}.
 \label{eq:comp-real-combined-kernel}
\end{align}
All these coefficients follow from the preceding rational partial
fractions and inner boundary differences.

Consider the polynomial
\begin{equation}
 Q_b(u)=u^2+2u+s=(u+1)^2+s-1.
\end{equation}
For $s>1$ it is irreducible over the real numbers. Its two roots
$b_\pm=-1\pm\ii\sqrt{s-1}$ would introduce an additional pair of GPL
letters if the sectors were integrated separately. Polynomial division
shows that their entire pole family cancels. To specify this statement
precisely, let $\nu_b(K)$ be the multiplicity of $Q_b$ in the reduced
denominator of a rational coefficient $K$. The inner integrals give
\begin{equation}
\begin{array}{c|ccc}
 j&\nu_b(K_{R_0,j})&\nu_b(K_{D,j})&
                 \nu_b(K_{R_0,j}+K_{D,j})\\ \hline
 0&2&2&0\\
 1&0&0&0\\
 2&3&3&0\\
 3&3&3&0
\end{array}.
 \label{eq:comp-real-b-pole-orders}
\end{equation}
The cancellation can be checked without factoring $Q_b$ into its roots.
For each row, put the two coefficients over a common denominator
$Q_b^n d$, where $d$ has no factor $Q_b$, and write their numerator
sum as $N$. Euclidean division gives
\begin{equation}
 N=Q_b^n M+R,\qquad \deg_uR<2n,\qquad R=0.
\end{equation}
Cancelling $Q_b^n$ then gives $M/d$. This is an exact identity over
the rational functions of $H=s-1$. The energy-logarithm and single-$W$
coefficients have no $Q_b$ factor, so none is reintroduced in
\eqref{eq:comp-real-combined-kernel}.

The preferred outer integration applies partial fractions to these
combined coefficients. Products of logarithms are converted to GPL
words by the shuffle identity, and repeated poles are reduced by the
integration-by-parts operations already given. The letters $b_\pm$
need never be introduced, and their rational pole terms and six
weight-two functions need never be evaluated separately. This reduction
also leaves no boundary constant to determine for the removed group:
its integrand vanishes identically at every common finite cutoff.
Linearity then preserves the identity when the endpoint finite parts
and their previously derived cutoff conversions are taken. The exact
coefficient construction and the polynomial divisions are implemented
in \nolinkurl{compton_real_combined_kernel.py} using only the inner
rational kernels.
